% !PRINTED 420
\weberpage{420}
Nun haben wir im \S.\,56 zwei Formen $f$ und $H$ vom $6^{\mathrm{ten}}$ und $8^{\mathrm{ten}}$ Grade kennen gelernt, die durch die Tetraëdersubstitutionen, wenn sie homogen und mit der Determinante 1 genommen werden, absolut ungeändert bleiben. Diese Formen können wir freilich nicht geradezu in der Form anwenden, wie sie dort gegeben sind, weil dort die Tetraëdergruppe in anderer Gestalt vorausgesetzt war. Wir finden aber die erste dieser Formen sehr einfach, wenn wir uns erinnern, dass ihre Wurzeln die zweizähligen Pole der Tetraëdergruppe waren, also die Wurzeln der drei Gleichungen
\[
x = \psi(x), \quad x = \chi(x), \quad x = \psi\chi(x)
\]
oder
\[
x^2 + 1 = 0, \quad x^2 - 2\omega x - 1 = 0, \quad \omega x^2 + 2x - \omega = 0,
\]
und man erhält also mit Benutzung der Relation
\[
\omega^2 + \omega = 1
\]
die Gleichung für die zweizähligen Pole
\[
(x^2+1)(x^4+2x^3-6x^2-2x+1) = 0.
\]

Daher wird die gesuchte Tetraëderform $6^{\mathrm{ten}}$ Grades in den homogenen Variablen $x_1, x_2$
\weq{5}{t(x_1,x_2) = x_1^6 + 2x_1^5x_2 - 5x_1^4x_2^2 - 5x_1^2x_2^4 - 2x_1x_2^5 + x_2^6.}

Die zweite Tetraëderform, die wir suchen, ist die \gesperrt{Hesse}'sche Determinante hiervon.

Setzen wir
\[
\tilde\omega(x_1,x_2) = \dfrac{1}{5^2 \cdot 4^2} \left\{ t''(x_1,x_1)\, t''(x_2,x_2) - [t''(x_1,x_2)]^2 \right\},
\]
so ergiebt eine einfache Rechnung
\weq{6}{\tilde\omega(x_1,x_2) = -(x_1^8+x_2^8) + (x_1^7x_2 - x_1x_2^7) - 7(x_1^6x_2^2+x_1^2x_2^6) - 7(x_1^5x_2^3-x_1^3x_2^5).}

Wir führen noch die drei Ikosaëderformen $12^{\mathrm{ten}}$, $20^{\mathrm{sten}}$ und $30^{\mathrm{sten}}$ Grades an, die ja auch durch die Tetraëdersubstitutionen ungeändert bleiben (\S.\,60):
\weq{7}{f(x_1,x_2) = x_1x_2(x_1^{10}+11x_1^5x_2^5-x_2^{10})}
\weq{8}{H(x_1,x_2) = -(x_1^{20}+x_2^{20}) + 228(x_1^{15}x_2^5-x_1^5x_2^{15}) - 494x_1^{10}x_2^{10}}
\weq{9}{T(x_1,x_2) = (x_1^{30}+x_2^{30}) + 522(x_1^{25}x_2^5-x_1^5x_2^{25}) - 10005(x_1^{20}x_2^{10}+x_2^{20}x_1^{10}).}
